\documentclass[a4paper]{arrowhead}

\usepackage[yyyymmdd]{datetime}
\usepackage{etoolbox}
\usepackage[utf8]{inputenc}
\usepackage[T1]{fontenc}
\usepackage{multirow}
\usepackage{longtable}
\usepackage{tikz}
\usetikzlibrary{positioning,shapes.multipart,arrows.meta}

\renewcommand{\dateseparator}{-}
\setlength{\emergencystretch}{3em}
\lstset{language=HTTP,breakatwhitespace=false,columns=fullflexible}

%% Special references
\newcommand{\mref}[2][]{{\textcolor{ArrowheadPurple}{\hyperref[sec:model:#2]{\ifstrempty{#1}{#2}{#1}}}}}
\newcommand{\pref}[1]{{\textcolor{ArrowheadGrey}{\hyperref[sec:model:primitives:#1]{#1}}}}
%%

\begin{document}
\sloppy

%% Arrowhead Document Properties
\ArrowheadTitle{Blacklist}
\ArrowheadType{System Description}
\ArrowheadTypeShort{SysD}
\ArrowheadVersion{5.2.2}
\ArrowheadDate{2026-10-30}
\ArrowheadAuthor{Jerker Delsing}
\ArrowheadStatus{RELEASE CANDIDATE}
\ArrowheadContact{jerker.delsing@ltu.se}
\ArrowheadFooter{\href{www.arrowhead.eu}{www.arrowhead.eu}}
\ArrowheadSetup
%%

%% Front Page
\begin{center}
  \vspace*{1cm}
  \huge{\arrowtitle}

  \vspace*{0.2cm}
  \LARGE{\arrowtype}
  \vspace*{1cm}

  \vspace*{\fill}

  % Front Page Image

  \vspace*{1cm}
  \vspace*{\fill}

  % Front Page Abstract
  \begin{abstract}
    This document describes the Blacklist support system of the Eclipse Arrowhead 5.2.2 framework, as implemented in \texttt{ah5-Blacklist-java-spring} of the primary Java source tree (\texttt{Eclipse-Arrowhead-java}) and deployed with the standalone Compose definition of that folder.
  \end{abstract}

  \vspace*{1cm}
\end{center}

\newpage
%%

%% Table of Contents
\tableofcontents
\newpage
%%

\section{Overview}
\label{sec:overview}

This document describes the Blacklist system. The Blacklist keeps the list of systems that are excluded from the local cloud. Every system with the blacklist filter enabled asks it before serving a request, and Dynamic\allowbreak{}Service\allowbreak{}Orchestration removes blacklisted providers from its results.

\begin{figure}[ht!]
  \centering
  \begin{tikzpicture}[font=\small, >=Stealth]
    \node[draw=ArrowheadBlue, very thick, minimum width=4.2cm, minimum height=2.2cm, align=center, fill=ArrowheadBlue!10] (s)
      {\guillemotleft{}System\guillemotright{}\\\textbf{Blacklist}};
    \node[draw=ArrowheadTeal, rectangle, align=left, left=1.3cm of s, text width=4.6cm, inner sep=5pt] (p)
      {\textbf{Produced services}\\ \texttt{blacklist\allowbreak{}Management}\\ \texttt{blacklistDiscovery}\\ \texttt{generalManagement}\\ \texttt{sysMonitor}};
    \node[draw=ArrowheadOrange, rectangle, align=left, right=1.3cm of s, text width=4.6cm, inner sep=5pt] (c)
      {\textbf{Consumed}\\ AMQP broker (SysMonitor)\\ Authentication\\ ConsumerAuthorization\\ ServiceRegistry};
    \draw[->, thick, ArrowheadTeal] (s) -- (p);
    \draw[->, thick, ArrowheadOrange] (s) -- (c);
  \end{tikzpicture}
  \caption{SysML-style block view of Blacklist with its produced and consumed services.}
  \label{fig:system}
\end{figure}

The rest of this document is organized as follows.
In Section \ref{sec:prior_art}, we reference major prior art capabilities of the system.
In Section \ref{sec:use}, we describe the intended usage of the system.
In Section \ref{sec:properties}, we describe fundamental properties provided by the system.
In Section \ref{sec:delimitations}, we describe delimitations of capabilities of the system.
In Section \ref{sec:services}, we describe the abstract service functions consumed or produced by the system.
In Section \ref{sec:security}, we describe the security capabilities of the system.

\newpage

\subsection{Significant Prior Art}
\label{sec:prior_art}

New in Arrowhead 5 (AITIA 5.2.2 implementation).

\subsection{How This System Is Meant to Be Used}
\label{sec:use}

Operators or the onboarding procedure add a system with \texttt{blacklist\allowbreak{}Management.\allowbreak{}create}; from then on every system with \texttt{enable.\allowbreak{}blacklist.\allowbreak{}filter=\allowbreak{}true} refuses its requests (403) and orchestration no longer returns it. \texttt{remove} lifts the entry.

\subsection{System functionalities and properties}
\label{sec:properties}

\subsubsection{Functional properties of the system}

\begin{itemize}
  \item Store blacklist entries with reason, creator and optional expiry; inactivate them instead of deleting.
  \item Answer \texttt{check} for a system name and \texttt{lookup} for the requester's own entries.
  \item Expire entries at \texttt{expiresAt}.
\end{itemize}

\subsubsection{Configuration of system properties}

The system is configured with Spring properties (\texttt{application.properties}, overridden by environment variables in the container deployment). The properties that shape its behaviour are:

\noindent\begin{longtable}{| p{5.4cm} | p{3.4cm} | p{6.4cm} |} \hline
\rowcolor{gray!33} Property & Default & Meaning \\ \hline \endhead
\texttt{server.\allowbreak{}port} & 8464 & HTTPS port. \\ \hline
\texttt{authentication.\allowbreak{}policy} & declared & Standalone deployment: \texttt{outsourced} with Authentication identity tokens and mutual TLS. \\ \hline
\texttt{management.\allowbreak{}policy} / \texttt{management.\allowbreak{}whitelist} & authorization & Access to the management services. \\ \hline
\texttt{normalization.\allowbreak{}mode} & extended & Name normalization. \\ \hline
\texttt{mqtt.\allowbreak{}api.\allowbreak{}enabled} & false & Offer the MQTT interfaces. \\ \hline
\texttt{sysmonitor.\allowbreak{}enabled,\allowbreak{} sysmonitor.\allowbreak{}broker.\allowbreak{}*,\allowbreak{} sysmonitor.\allowbreak{}ssl.\allowbreak{}*} & true, amqp-broker:5671, server.ssl stores & Native \texttt{sysMonitor} publisher of the common library. \\ \hline
\texttt{arrowhead.\allowbreak{}ssl.\allowbreak{}crl-\allowbreak{}file} & (unset) & CRL checked for every client certificate (standalone deployment: \texttt{/\allowbreak{}crl/\allowbreak{}ca.\allowbreak{}crl}). \\ \hline
\end{longtable}

Secrets (database, keystore and identity passwords, secret keys) have no usable defaults; the standalone deployment reads them from per-deployment Compose file secrets under \texttt{/\allowbreak{}run/\allowbreak{}secrets}.

The standalone deployment of this folder uses \texttt{config/\allowbreak{}application.\allowbreak{}properties} (copied into the image) with \texttt{server.\allowbreak{}port}=8464, \texttt{authentication.\allowbreak{}policy}=\texttt{outsourced}, \texttt{service.\allowbreak{}interface.\allowbreak{}policy}=\texttt{CERT\_AUTH}, \texttt{management.\allowbreak{}policy}=\texttt{sysop-\allowbreak{}only}, \texttt{service.\allowbreak{}authorization.\allowbreak{}enabled}=\texttt{true}, \texttt{mqtt.\allowbreak{}api.\allowbreak{}enabled}=\texttt{false}, \texttt{arrowhead.\allowbreak{}ssl.\allowbreak{}crl-\allowbreak{}file}=\texttt{/\allowbreak{}crl/\allowbreak{}ca.\allowbreak{}crl}; secrets are read from Compose file secrets under \texttt{/\allowbreak{}run/\allowbreak{}secrets} (\texttt{spring.\allowbreak{}config.\allowbreak{}import=\allowbreak{}configtree:\allowbreak{}/\allowbreak{}run/\allowbreak{}secrets/}). \texttt{enable.\allowbreak{}blacklist.\allowbreak{}filter} is not enabled there, so incoming requests are not checked against the Blacklist; blacklisted providers are still removed by Dynamic\allowbreak{}Service\allowbreak{}Orchestration.

\subsubsection{Data stored by the system}

Blacklist entries and log entries (database \texttt{ah\_\allowbreak{}blacklist}).

\subsubsection{Non functional properties}

Called on the request path of every filtering system; responses are small database lookups.

\subsubsection{Stateful or stateless}

Stateful: entries are persistent.

\subsection{Important Delimitations}
\label{sec:delimitations}

The Blacklist does not enforce anything itself. Systems in \texttt{blacklist.\allowbreak{}check.\allowbreak{}exclude.\allowbreak{}list} (by default the core systems) are never checked.

\newpage

\section{Services}
\label{sec:services}

This section describes the produced and consumed services. Each produced service is specified by an SD document and one IDD document per interface.

\subsection{Produced services}

\noindent\begin{longtable}{| p{5.3cm} | p{1.3cm} | p{2.6cm} | p{1.6cm} | p{3.6cm} |} \hline
\rowcolor{gray!33} Service & Version & Interfaces & SD & IDD \\ \hline \endhead
\texttt{blacklist\allowbreak{}Management} & 1.0.0 & HTTP(S), MQTT(S) & \cite{bl-blacklistManagement-sd} & \cite{bl-blacklistManagement-http-idd}, \cite{bl-blacklistManagement-mqtt-idd} \\ \hline
\texttt{blacklistDiscovery} & 1.0.0 & HTTP(S), MQTT(S) & \cite{bl-blacklistDiscovery-sd} & \cite{bl-blacklistDiscovery-http-idd}, \cite{bl-blacklistDiscovery-mqtt-idd} \\ \hline
\texttt{generalManagement} & 1.0.0 & HTTP(S), MQTT(S) & \cite{bl-generalManagement-sd} & \cite{bl-generalManagement-http-idd}, \cite{bl-generalManagement-mqtt-idd} \\ \hline
\texttt{sysMonitor} & 1.0.0 & AMQP(S) & \cite{bl-sysMonitor-sd} & \cite{bl-sysMonitor-amqp-idd} \\ \hline
\end{longtable}

The HTTP interfaces are always offered, the MQTT interfaces with \texttt{mqtt.api.enabled=true} and the AMQP \texttt{sysMonitor} interface with \texttt{sysmonitor.enabled=true}. The unregistered endpoints of the implementation are: \texttt{GET /\allowbreak{}blacklist/\allowbreak{}monitor/\allowbreak{}echo}.

\subsection{Consumed services}

\noindent\begin{longtable}{| >{\raggedright\arraybackslash}p{2.9cm} | >{\raggedright\arraybackslash}p{3.1cm} | >{\raggedright\arraybackslash}p{2.5cm} | p{3.6cm} | >{\raggedright\arraybackslash}p{2.9cm} |} \hline
\rowcolor{gray!33} Provider & Service & Operations & Use & Condition \\ \hline \endhead
ServiceRegistry & systemDiscovery, serviceDiscovery & register, revoke, lookup & Registration of the system and its services; lookup of the bootstrap services. & always \\ \hline
Authentication & identity & identity-login, identity-verify & Own identity token; verification of requester tokens. & \texttt{authentication.\allowbreak{}policy=\allowbreak{}outsourced} \\ \hline
Consumer\allowbreak{}Authorization & authorization & verify & Management filter decisions. & \texttt{management.\allowbreak{}policy=\allowbreak{}authorization} \\ \hline
AMQP broker (SysMonitor) & sysMonitor (publisher) & echo, power, status & Own observations. & \texttt{sysmonitor.\allowbreak{}enabled=\allowbreak{}true} \\ \hline
\end{longtable}

Except for the bootstrap services, which are looked up in the ServiceRegistry, the system reaches other systems' services through DynamicServiceOrchestration (\texttt{http.collector.mode=ORCHESTRATION}, results cached for 300 s); every such call needs a ConsumerAuthorization grant.

\newpage

\section{Security}
\label{sec:security}

The system can run in an insecure mode (\texttt{server.ssl.enabled=false}, plain HTTP, identity declared by the requester) for development, and in the Arrowhead secure mode with HTTPS and X.509 certificates following the Arrowhead certificate profile (\texttt{<System>.<Cloud>.<Organization>.arrowhead.eu}). Certificates not issued under the cloud CA are refused by the TLS handshake. The standalone deployment of this folder (\texttt{config/\allowbreak{}application.\allowbreak{}properties} and \texttt{compose.yaml}) runs the secure mode.

\subsection{Security model}


\begin{itemize}
  \item \textbf{Protocols:} HTTPS (TLS 1.3 and 1.2) with mandatory client certificates; MQTT over TLS; AMQP over TLS for \texttt{sysMonitor}.
  \item \textbf{Data protection:} transport encryption by TLS; the database connection uses TLS with server verification in the standalone deployment. Secrets are kept out of the container environment (Compose file secrets).
  \item \textbf{System authentication:} by the authentication policy: \texttt{certificate}, \texttt{outsourced} (identity token verified with Authentication) or \texttt{declared} (insecure); the standalone deployment uses \texttt{outsourced}. An identity token is accepted only together with the client certificate of the system it was issued to (token-to-certificate binding), and every client certificate is checked against the cloud CA's CRL (fail closed).
  \item \textbf{Authorization of produced services:} management services are protected by the management filter (\texttt{management.policy}: \texttt{SYSOP\_ONLY}, \texttt{WHITELIST} or \texttt{AUTHORIZATION}); with \texttt{service.authorization.enabled=true} ConsumerAuthorization decides each call; an optional blacklist filter refuses blacklisted requesters. Sysop identities bypass the management filter.
  \item Off-boarded systems are blacklisted by Onboarding so that their remaining credentials are refused.
\end{itemize}

For the Arrowhead certificate profile see github.com/eclipse-arrowhead/documentation.


\bibliographystyle{IEEEtran}
\bibliography{bibliography}

\newpage

\section{Revision History}
\subsection{Amendments}

\noindent\begin{tabularx}{\textwidth}{| p{1cm} | p{3cm} | p{2cm} | X | p{4cm} |} \hline
\rowcolor{gray!33} No. & Date & Version & Subject of Amendments & Author \\ \hline

1 & 2026-10-22 & \arrowversion & Document created from the Arrowhead LaTeX templates for the 5.2.2 Java Spring implementation, including the local extensions of this workspace. & Jerker Delsing \\ \hline
2 & 2026-10-30 & \arrowversion & Checked against the primary Java source tree (Eclipse-Arrowhead-java): standalone deployment, configuration and security statements updated. & Claude Opus 5.5 (draft) \\ \hline
\end{tabularx}

\subsection{Quality Assurance}

\noindent\begin{tabularx}{\textwidth}{| p{1cm} | p{3cm} | p{2cm} | X |} \hline
\rowcolor{gray!33} No. & Date & Version & Approved by \\ \hline

1 & & \arrowversion  &  \\ \hline

\end{tabularx}

\end{document}
